\documentclass[11pt,a4paper]{article}
\usepackage[T1]{fontenc}
\usepackage[utf8]{inputenc}
\usepackage{lmodern}
\usepackage{amsmath,amssymb,amsthm,mathtools,mathrsfs}
\usepackage{graphicx,float,xcolor,multicol}
\usepackage[shortlabels]{enumitem}
\usepackage[margin=27mm]{geometry}
\usepackage{microtype,needspace,placeins}
\usepackage{tikz,tikz-cd}
\usetikzlibrary{arrows.meta,calc,positioning,patterns,decorations.pathreplacing}
\usepackage[colorlinks=true,linkcolor=teal,urlcolor=teal,citecolor=teal]{hyperref}
\setlength{\parindent}{0pt}
\setlength{\parskip}{.6em}
\setcounter{secnumdepth}{0}
\allowdisplaybreaks
\raggedbottom
\providecommand{\cross}{\times}
\DeclareUnicodeCharacter{2020}{\ensuremath{\dagger}}
\DeclareMathOperator{\supp}{supp}
\DeclareMathOperator*{\esssup}{ess\,sup}
\providecommand{\chapter}{\section}
\newenvironment{tool}[2]{\subsection*{Tool #1 --- #2}}{}
\newcommand{\ToolRef}[1]{Tool #1}
\newcommand{\FigureTag}[2]{\textit{Figure #1}}
\newcommand{\Result}[1]{\subsection*{#1}}
\newcommand{\Application}[1]{\subsubsection*{Application\if\relax\detokenize{#1}\relax\else\space--- #1\fi}}
\newcommand{\ProofHeading}{\par\textit{Proof.}\space}
\newcommand{\ProofOf}[1]{\par\textit{Proof of #1.}\space}
\newcommand{\RemarkHeading}{\par\textbf{Remark.}\space}
\newcommand{\NoteHeading}{\par\textbf{Note.}\space}
\newcommand{\ExerciseHeading}{\par\textbf{Exercise.}\space}

% Rendering support only; mathematical text is in each independent post.
\usepackage{booktabs,array,longtable}
\definecolor{HandoutBlue}{HTML}{2F6690}
\definecolor{HandoutNavy}{HTML}{17365D}
\newtheorem{theorem}{Theorem}
\newtheorem{lemma}[theorem]{Lemma}
\newtheorem{proposition}[theorem]{Proposition}
\newtheorem{corollary}[theorem]{Corollary}
\newtheorem{conjecture}[theorem]{Conjecture}
\newtheorem{definition}{Definition}
\newtheorem{example}{Example}
\newtheorem{namedtoolcounter}{Tool}
\newenvironment{namedtool}[1]{\begin{namedtoolcounter}[#1]}{\end{namedtoolcounter}}
\newtheorem*{claim}{Claim}
\newtheorem*{remark}{Remark}
\newtheorem*{exercise}{Exercise}
\newtheorem*{question}{Question}
\newtheorem*{observation}{Observation}
\newcommand{\SourceChapter}[2]{\section*{#2}}
\newcommand{\Lecture}[2]{\section*{Lecture #1: #2}}
\newcommand{\unclear}[1]{\textit{[unclear: #1]}}
\newcommand{\sourcegap}{\textit{[unfinished in the source]}}
\DeclareMathOperator{\dist}{dist}
\DeclareMathOperator{\id}{id}
\DeclareMathOperator{\Hol}{Hol}
\DeclareMathOperator{\Per}{Per}
\DeclareMathOperator{\Fix}{Fix}
\DeclareMathOperator{\diam}{diam}
\DeclareMathOperator{\Var}{Var}
\DeclareMathOperator{\Leb}{Leb}
\DeclareMathOperator{\Hom}{Hom}
\DeclareMathOperator{\End}{End}
\DeclareMathOperator{\Ann}{Ann}
\DeclareMathOperator{\Spec}{Spec}
\DeclareMathOperator{\Max}{Max}
\DeclareMathOperator{\Ob}{Ob}
\DeclareMathOperator{\Tor}{Tor}
\DeclareMathOperator{\Ass}{Ass}
\DeclareMathOperator{\Mor}{Mor}
\DeclareMathOperator{\Fun}{Fun}
\DeclareMathOperator{\Nat}{Nat}
\DeclareMathOperator{\Cone}{Cone}
\DeclareMathOperator{\MCG}{MCG}
\DeclareMathOperator{\Aut}{Aut}
\DeclareMathOperator{\Deck}{Deck}
\DeclareMathOperator{\SL}{SL}
\DeclareMathOperator{\PSL}{PSL}
\DeclareMathOperator{\Area}{Area}
\DeclareMathOperator{\im}{im}
\DeclareMathOperator{\PSh}{PSh}
\newcommand{\CRing}{\mathsf{CRings}}
\newcommand{\AMod}{A\text{-}\mathsf{Mod}}
\newcommand{\Ab}{\mathsf{Ab}}
\newcommand{\Z}{\mathbb Z}
\newcommand{\Q}{\mathbb Q}
\newcommand{\R}{\mathbb R}
\newcommand{\C}{\mathbb C}
\newcommand{\N}{\mathbb N}
\newcommand{\kk}{\mathbb k}
\renewcommand{\aa}{\mathfrak a}
\newcommand{\bb}{\mathfrak b}
\newcommand{\pp}{\mathfrak p}
\newcommand{\qq}{\mathfrak q}
\newcommand{\mm}{\mathfrak m}
\newcommand{\nn}{\mathfrak n}
\newcommand{\rad}{\mathrm r}
\newcommand{\cat}[1]{\mathsf{#1}}
\setcounter{secnumdepth}{0}
\allowdisplaybreaks

\graphicspath{{figures/lemma-book/}}
\begin{document}
\section*{Congruences, Diophantine Equations, and Infinite Descent}
% BLOG-CONTENT-BEGIN
\setcounter{theorem}{45}\begin{lemma}[Two obvious congruences]
$x\equiv x\pmod n$ and $x\equiv0\pmod x$ have surprisingly ingenious applications.
\end{lemma}
Find the largest positive integer $n$ such that $n+10\mid n^3+100$.
Since $n+10\equiv0\pmod{n+10}$, we have $n\equiv-10\pmod{n+10}$.
Thus $n^3+100\equiv-900\equiv0\pmod{n+10}$, so $n+10\mid900$.
The maximum value is $n=890$.

\setcounter{theorem}{46}\begin{lemma}[Lagrange's identity]
$(a^2+b^2)(c^2+d^2)=(ac\pm bd)^2+(ad\mp bc)^2$.
\end{lemma}
Represent $m^6+n^6$ as a sum of two squares other than $(m^3)^2+(n^3)^2$.
\[
m^6+n^6=(m^2)^3+(n^2)^3
=(m^2+n^2)(m^4-m^2n^2+n^4)
=(m^2+n^2)\bigl((m^2-n^2)^2+(mn)^2\bigr).
\]
This suggests using Lagrange's identity. Set $a=m$, $b=n$, $c=m^2-n^2$, and $d=mn$.
Then $m^6+n^6=(m^3-2mn^2)^2+(n^3-2nm^2)^2$.

\setcounter{theorem}{47}\begin{lemma}[A prime in one denominator]
If $\sum_i a_i/b_i=N/D$ is in lowest terms and exactly one denominator $b_k$ has maximal positive $p$-adic valuation, with $p\nmid a_k$, then $p\mid D$ and $p\nmid N$.
\end{lemma}
Prove that $1+1/2+1/3+\cdots+1/n$ is never an integer for $n>1$.
Let $2^r$ be the largest power of $2$ not exceeding $n$, and let $L=\operatorname{lcm}(1,2,\ldots,n)$.
In $\sum_{k=1}^n L/k$, the term $L/2^r$ is odd and every other term is even.
Thus the numerator is odd before cancellation while $L$ is even, so the reduced denominator is even.

% Source page 23 - left
\setcounter{theorem}{48}\begin{lemma}[Products of perfect powers]
If $a_i$ is a perfect $n$th power for every $1\leq i\leq k$, then $\prod_{i=1}^k a_i$ is a perfect $n$th power.
\end{lemma}
For $n\geq0$, prove that $n+2$ and $n^2+n+1$ cannot both be perfect cubes.
If they were, $(n+2)(n^2+n+1)=a^3$ for some $a$. But
\[
n^3+n^2+n+2n^2+2n+2=a^3
\Longrightarrow(n+1)^3+1=a^3,
\]
a contradiction.

\setcounter{theorem}{50}\begin{lemma}[Fermat's last theorem]
$x^n+y^n=z^n$ has no solution for $n>2$.
\end{lemma}
Fermat's last theorem shortens the following argument.
If $x,y,z$ are positive natural numbers and $n>1$, with $x^n+y^n=z^n$, show that $\min(x,y,z)>n$.
Fermat's last theorem gives $n=2$, since $n>1$.
Thus $x^2+y^2=z^2$. If, for example, $x\leq2$, then
$(z-y)(z+y)=x^2\leq4$, which has no solution with $y>0$; the same argument applies to $y$, while $z>x,y$.

\emph{Experimentation, observation, and insight bottled together possess
the capability to perform wonders.}
Given
\[
\overline{abc}_{10}+\overline{acb}_{10}+\overline{bac}_{10}
+\overline{bca}_{10}+\overline{cab}_{10}+\overline{cba}_{10}=3194,
\]
find $\overline{abc}_{10}$.
This is a problem from my fruitful problem-solving experience. I employed
a unique, original technique. Let $\overline{abc}_{10}=100a+10b+c$.
The notes write
\[
3194=122a+212b+221c=f(a,b,c).
\]
Guess: $f(7,6,8)=3894$. This exceeds $3194$ by $700$.
Since $f(4,1,0)=700$, we obtain
$f(7-4,6-1,8-0)=f(3,5,8)=3194$.
Thus $\overline{abc}=358$.

In fact, every digit occurs twice in each
place among the six permutations, so the left-hand side is
$222(a+b+c)$. Since $222\nmid3194$, the problem as written has no
solution. Thus the coefficients $122,212,221$ used in the proposed
calculation do not match the displayed sum of all six permutations; the
numerical constant or the displayed list of permutations must contain a
typographical error.

\subsection{Ramanujan's ingenious idea}
\setcounter{theorem}{53}\begin{lemma}[Adjusting a sum with ones]
To arrange
\[
A+B+C+D+\cdots=P+Q+R+S+\cdots,\qquad
A^AB^BC^CD^D\cdots=P^PQ^QR^RS^S\cdots,
\]
first find quantities satisfying the product equation, then multiply by enough $1^1$ terms to satisfy the sum equation.
For example, $2^2\cdot6^6=3^3\cdot3^3\cdot4^4$ becomes
$1^1\cdot1^1\cdot2^2\cdot6^6=3^3\cdot3^3\cdot4^4$,
and $1+1+2+6=3+3+4$.
\end{lemma}
\paragraph{Application: APMO (2014).}
For a positive integer $m$, let $S(m)$ and $P(m)$ denote its digit sum and digit product.
Show that, for each positive $n$, there are natural numbers $a_1,\ldots,a_n$ such that
\[
S(a_1)<S(a_2)<\cdots<S(a_n),\qquad S(a_i)=P(a_{i+1})\quad(1\leq i<n).
\]
After some analysis, I would stick to the sequence
$(a_n)=(11,\allowbreak211,\allowbreak221111,\allowbreak22211111111,\allowbreak\ldots)$.
Guessing the pattern is not difficult. This is somewhat close to a magical solution!

\setcounter{theorem}{55}\begin{lemma}[Subtract the two sides]
To show $A(x)=B(x)$, define $f(x)=A(x)-B(x)$ and try to prove $f(x)=0$ everywhere.
\end{lemma}
\paragraph{Application: Hermite's identity.}
% Source page 25 - right
Prove that
$\lfloor nx\rfloor=\lfloor x\rfloor+\lfloor x+1/n\rfloor+\cdots+\lfloor x+(n-1)/n\rfloor$.
Define $f(x)=\sum_{k=0}^{n-1}\lfloor x+k/n\rfloor-\lfloor nx\rfloor$.
Then
\[
f(x+1/n)=\lfloor x+1/n\rfloor+\cdots+\lfloor x+(n-1)/n\rfloor
+\lfloor x+1\rfloor-\lfloor nx+1\rfloor.
\]
Since $\lfloor x+k\rfloor=\lfloor x\rfloor+k$ for integer $k$, we get $f(x+1/n)=f(x)$.
Thus $f$ has period $1/n$. For $0\leq x\leq1/n$, $f(x)=0$; therefore it vanishes everywhere, proving the identity.

\setcounter{theorem}{63}\begin{lemma}[Floor and fractional part]
$\lfloor x\rfloor=x-\{x\}$, where $\{x\}\in[0,1)$.
Thus $0\leq\{x\}<1$ and $x-1<\lfloor x\rfloor\leq x$.
\end{lemma}
This useful lemma can turn a floor-function equation into an elementary inequality problem.
\paragraph{Application.}
Find all real numbers satisfying
$\lfloor x/2\rfloor+\lfloor x/4\rfloor+\lfloor x/8\rfloor=2014$.
Write $x/8=q+u$, where $q\in\mathbb Z$ and $0\leq u<1$. Then
\[
\left\lfloor\frac x2\right\rfloor+
\left\lfloor\frac x4\right\rfloor+
\left\lfloor\frac x8\right\rfloor
=7q+\lfloor4u\rfloor+\lfloor2u\rfloor.
\]
The last two terms have sum $0,1,3$, or $4$, whereas $2014\equiv5\pmod7$.
Therefore the equation has no real solutions.

Here are two problems with particularly striking solutions.
Now let me present a problem and its standout solution!
Find $a,b,c>0$ such that $1/2014=1/a+1/b+1/c$.
This was asked in class. One student boldly gave the unexpected solution:
``We know $1=1/2+1/3+1/6$; now divide this by $2014$ to obtain $a,b,c$!''

\setcounter{theorem}{66}\begin{lemma}[Transform one solution into another]
In some Diophantine equations, it helps to observe that if $(a,b,c)$ is a solution, then $(f(a),f(b),f(c))$ is also a solution.
\end{lemma}
\paragraph{Application: a classical IMO problem.}
Let $a,b$ be positive integers with $ab+1\mid a^2+b^2$.
Show that $(a^2+b^2)/(ab+1)$ is an integer square.
Write $(a^2+b^2)/(ab+1)=k$, or $a^2-kab+b^2=k$.
If $(a,b)$ solves this equation, then a direct substitution shows that $(b,kb-a)$ also solves it.
If $a=b$, then $(2-k)a^2=k>0$, so $k=1$ is a solution.
Without loss of generality, assume $a>b$. Then
\[
k=\frac{a^2+b^2}{ab+1}<\frac{a^2+ab}{ab+1}<\frac{a^2+ab}{ab}=\frac ab+1.
\]
Choose a solution $(a,b)$ with $b$ as small as possible among all positive solutions for this value of $k$.
Put $c=kb-a$. The quadratic equation gives $ac=b^2-k$.
If $c<0$, then $a-kb>0$ and $k-b^2=a(a-kb)\geq a$, contradicting $k<a/b+1\leq a+1$.
Thus $c\geq0$, and the above bound gives $c<b$.
If $c>0$, then $(b,c)$ is a smaller positive solution, contradicting the choice of $b$.
Therefore $c=0$, and $k=b^2$ is an integer square.

Here is a useful parametric construction.
Prove that $x^2=y^3+z^5$ has infinitely many positive integer solutions.
The following family gives the required solutions.
% Source page 30 - right
\[
x_n=n^{10}(n+1)^8,\qquad y_n=n^7(n+1)^5,\qquad z_n=n^4(n+1)^3,
\qquad n\in\mathbb Z_+.
\]

\subsection{Diophantine equations and normalization}
\setcounter{theorem}{79}\begin{lemma}[A greatest-common-divisor strategy]\label{lemma:lb1-82}
When solving some Diophantine equations in two variables, case-by-case analysis using $\gcd(x,y)=d$ turns out to be useful.
\end{lemma}
This strategy appears frequently in national olympiad problems.
\paragraph{Application: Belarus Mathematical Olympiad (2000).}
Find all pairs of positive integers $(m,n)$ satisfying
\[(m-n)^2(n^2-m)=4m^2n.\]
Let $\gcd(m,n)=d$, so $m=dx$ and $n=dy$, where $\gcd(x,y)=1$.
Then
\[(dx-dy)^2(d^2y^2-dx)=4d^3x^2y
\quad\Longrightarrow\quad (x-y)^2(dy^2-x)=4x^2y.\]
Reducing modulo $y$ gives $y\mid x^3$. Since $\gcd(x,y)=1$, we have $y=1$.
Hence $(x-1)^2(d-x)=4x^2$. Since $x\neq1$ and $\gcd(x,x-1)=1$, it follows that $(x-1)^2\mid4$.
Therefore $x=2$ or $3$, giving $d=18$ or $12$, respectively.
We end up with $(m,n)=(36,18)$ and $(36,12)$. Q.E.D.!

\setcounter{theorem}{87}\begin{lemma}[Dirichlet's theorem]\label{lemma:lb1-90}
An infinite arithmetic progression $a,a+b,a+2b,a+3b,\ldots$ contains infinitely many primes if $a,b$ are positive and $\gcd(a,b)=1$.
\end{lemma}
\paragraph{Application.}
Let $a_0=4$ and define $a_n=a_{n-1}^2-a_{n-1}$ for each positive integer $n$.
% Source page 50 - left
Prove that there are infinitely many primes that are factors of no term in the sequence.
Since $4^2\equiv4+2\pmod{10}$, we have $a_1\equiv2\pmod{10}$.
Now $2^2\equiv2+2\pmod{10}$, so $a_k\equiv2\pmod{10}$ for all $k\geq1$.
We have $a_1=12$. For $j\geq1$, write
\[
a_j=4t_j^2+6t_j+2,
\qquad t_1=1,
\qquad t_{j+1}=t_j(2t_j+3).
\]
A direct substitution verifies that this form is preserved by the recurrence. Moreover,
\[
\frac{a_{j+1}}{a_j}=a_j-1=4t_j^2+6t_j+1=(3t_j+1)^2-5t_j^2.
\]
The two integers $3t_j+1$ and $t_j$ are coprime. Hence, if a prime
$q>3$ divides $a_j-1$, then $5$ is a quadratic residue modulo $q$.
For every prime $q\equiv2$ or $3\pmod5$, quadratic reciprocity says that
$5$ is a nonresidue modulo $q$. Such a prime divides neither $a_1=12$
nor any later multiplier $a_j-1$, and therefore divides no term of the
sequence. Dirichlet's theorem supplies infinitely many such primes.

\subsection{Euler's formula and lifting indices}
\setcounter{theorem}{92}\begin{lemma}[Euler's formula]\label{lemma:lb1-95}
$a^{\phi(m)}\equiv1\pmod m$ if $\gcd(a,m)=1$.
\end{lemma}
\paragraph{Application: 20th IMO.}
Let $m,n$ be integers with $1\leq m<n$.
If $1978^m$ and $1978^n$ have the same last three digits, find the minimum of $m+n$.
We have $1978^n\equiv1978^m\pmod{1000}$, hence
$1978^m(1978^{n-m}-1)\equiv0\pmod{1000}$.
Thus $1978^m\equiv0\pmod8$ and $1978^{n-m}\equiv1\pmod{125}$.
The minimum of $m$ is $3$, so let $m=3$.
Now $1978^{n-3}\equiv1\pmod{125}$, and $\gcd(1978,125)=1$.
Euler's formula gives $1978^{\phi(125)}\equiv1\pmod{125}$, so the multiplicative order of $1978$ modulo $125$ divides $100$.
Also $1978^{n-3}\equiv3^{n-3}\pmod5$, and $1978^{n-3}\equiv1\pmod5$.
Hence $4\mid n-3$, since $3^4\equiv1\pmod5$.
The divisors of $100$ that are multiples of $4$ are $4,20$, and $100$.
Now
\[1978^4\equiv(-22)^4=4^2\cdot121^2\equiv(4\cdot(-4))^2
=(-16)^2=256\equiv6\not\equiv1\pmod{125}.\]
Also $1978^{20}\equiv26\not\equiv1\pmod{125}$.
Therefore the order is $100$, and the least possible value of $n-3$ is $100$, giving $\min n=103$.
Thus $\min(m+n)=\min m+\min n=103+3=106$.
Q.E.D.!
\Needspace{15\baselineskip}
\paragraph{Lifting the exponent.}\leavevmode\par
\setcounter{theorem}{93}\begin{lemma}[Lifting the Exponent]\label{lemma:lb1-96}
For a prime $p$, define $v_p(x)$ to be the exponent of the greatest power of $p$ that divides $x$.
If $v_p(x)=\alpha$, we also write $p^\alpha\mathbin{\|}x$.
Here are some lemmas:
\begin{enumerate}
\setcounter{enumi}{0}\item $v_p(xy)=v_p(x)+v_p(y)$, and $v_p(0)=\infty$.
\setcounter{enumi}{1}\item $v_p(x^n-y^n)=v_p(x-y)+v_p(n)$ if $p$ is odd, $p\mid x-y$, and $p\nmid xy$.
\setcounter{enumi}{2}\item $v_2(x^n-y^n)=v_2(x-y)+v_2(n)$ if $x,y$ are odd and $4\mid x-y$.
\setcounter{enumi}{3}\item $v_2(x^n-y^n)=v_2(x-y)+v_2(x+y)+v_2(n)-1$ if $n$ is even and $x,y$ are odd.
\setcounter{enumi}{4}\item $v_p(x^n+y^n)=v_p(x+y)+v_p(n)$ if $p$ is odd, $n$ is odd, $p\mid x+y$, and $p\nmid xy$.
\setcounter{enumi}{5}\item $v_p(x^n-y^n)=v_p(x-y)$ if $p$ is odd, $p\nmid nxy$, and $p\mid x-y$.
\setcounter{enumi}{6}\item $v_p(x^n+y^n)=v_p(x+y)$ if $p$ is odd, $n$ is odd, $p\nmid nxy$, and $p\mid x+y$.
\end{enumerate}
\end{lemma}
The lemma has many applications, provided its hypotheses are checked.
% Source page 57 - right
\paragraph{Application: Ireland (1996).}
Let $p$ be a prime and $a,n$ positive integers.
Prove that if $2^p+3^p=a^n$, then $n=1$.
If $p=2$, then $2^2+3^2=13$, so $n=1$.
Hence we may suppose that $p$ is an odd prime.
Now, if $p$ is odd, then $5\mid a$ and
$v_5(2^p+3^p)=v_5(2+3)+v_5(p)$, since $p$ is odd and $5\mid3+2$.
If $p=5$, we get $2^5+3^5=275=5^2\cdot11$, so again $n=1$.
Hence we may suppose that $p\neq5$.
Therefore $v_5(2^p+3^p)=1$.
Since $v_5(a^n)=n v_5(a)=1$, we obtain $n=1$.

\subsection{Fractional congruences}
\setcounter{theorem}{99}\begin{lemma}[Fractional congruences]\label{lemma:lb2-103}
If $a\equiv b\pmod n$ and $(b,n)=1$, then $a/b$ modulo $n$ means $ab^{-1}$ modulo $n$, and $a/b\equiv1\pmod n$.
\end{lemma}
\paragraph{Application: INMO practice test (2014).}
Given $1+1/2+\cdots+1/100=A/B$, with $(A,B)=1$, prove $5\nmid AB$.
Separate the sum as
\[
\sum_{\substack{1\leq k\leq100\\5\nmid k}}\frac1k
+\frac15\sum_{\substack{1\leq k\leq19\\5\nmid k}}\frac1k
+\frac1{25}\left(1+\frac12+\frac13+\frac14\right).
\]
The last term is $1/12$. Use the following fact:
\[
1+\frac12+\frac13+\frac14
\equiv\frac16+\frac17+\frac18+\frac19
\equiv\cdots\equiv\frac1{96}+\frac1{97}+\frac1{98}+\frac1{99}
\equiv1+2+3+4\pmod5.
\]
Thus $\sum_{5\nmid k,\,k\leq100}1/k\equiv200\equiv0\pmod5$,
so write this sum as $5X_1/Y_1$ with $(Y_1,5)=1$.
Further,
\[
1+\frac12+\frac13+\frac14=\frac{25}{12}\equiv0\pmod{25},\qquad
\frac16+\frac17+\frac18+\frac19\equiv21+18+22+14=75\equiv0\pmod{25}.
\]
Since $1/11\equiv-1/14$ and $1/12\equiv-1/13\pmod{25}$,
$1/11+1/12+1/13+1/14\equiv0\pmod{25}$.
Also
$1/16+1/17+1/18+1/19\equiv11+3+7+4\equiv0\pmod{25}$.
Thus $\sum_{5\nmid k,\,k\leq19}1/k=25X_2/Y_2$, where $(Y_2,5)=1$.
Consequently
\[
\frac AB=\frac{5X_1}{Y_1}+\frac15\frac{25X_2}{Y_2}+\frac1{12}
=\frac{60X_1Y_2+60X_2Y_1+Y_1Y_2}{12Y_1Y_2},
\]
so $5\nmid AB$.

The needed result is Zsigmondy's theorem.
\subsection{Zsigmondy's theorem}
\setcounter{theorem}{100}\begin{lemma}[Zsigmondy's theorem]\label{lemma:lb2-104}
If $a>b$ are positive integers with $\gcd(a,b)=1$ and $n>1$, then $a^n-b^n$ has a prime factor dividing none of $a^k-b^k$ for $k<n$,
with the exceptions $2^6-1^6$, and $n=2$ when $a+b$ is a power of $2$.
If $a>b$, $\gcd(a,b)=1$, and $n\geq2$, then
% Source page 8 - right
$a^n+b^n$ has a prime factor dividing none of $a^k+b^k$ for $k<n$, except $2^3+1^3$.
\end{lemma}
\paragraph{Application: IMO shortlist (2000).}
Find positive integer triples $(a,m,n)$ with $a^m+1\mid(a+1)^n$.
If $a=1$, every pair $(m,n)$ works, and if $m=1$, every pair $(a,n)$ works.
Suppose now that $a>1$ and $m>1$. Apart from $(a,m)=(2,3)$, Zsigmondy's theorem supplies a prime factor of $a^m+1$ not dividing $a+1$,
so it does not divide $(a+1)^n$. Finally, $2^3+1=3^2$ divides $3^n$ exactly when $n\geq2$.
Hence the remaining solutions are $(a,m,n)=(2,3,n)$ with $n\geq2$.

Solve the following using Zsigmondy.
\paragraph{Round 1: Fight. Japanese MO (2011).}
Find all positive integer quintuples $(a,n,p,q,r)$ such that
\[a^n-1=(a^p-1)(a^q-1)(a^r-1).\]
\paragraph{Round 2: Fight.}
Prove Fermat's last theorem, asserting that $x^n+y^n=z^n$ has no solutions
with $xyz\ne0$. Take $z$ to be prime!
\paragraph{Final round: Fight. Balkan MO (2009).}
Solve $5^x-3^y=z^2$ for $x,y,z\in\mathbb N$.

\paragraph{You used calculus in that!}
If $a^n+n$ is a multiple of $b^n+n$ for every positive integer $n$, prove
that $a=b$, where $a,b\in\mathbb N$.
Suppose there exists such an $n$ and seek $(a,b)$, with $a\geq b$.
Let $f(a)=a^n+n$ and $f(b)=b^n+n$, with $f(a)=kf(b)$,
$k\in\mathbb N$.
The proposed calculation differentiates:
\[
\frac{df(a)}{da}=\frac{d[kf(b)]}{db}
\Longrightarrow na^{n-1}=nkb^{n-1}
\Longrightarrow k=(a/b)^{n-1}.
\]
It then combines this with
$k=(a^n+n)/(b^n+n)=a^{n-1}/b^{n-1}$ to obtain
\[
n=\frac{a^{n-1}b^{n-1}(a-b)}{a^{n-1}-b^{n-1}},
\]
and announces a contradiction from Zsigmondy's theorem!
The differentiation step is not valid: $a$ and $b$ are fixed integers, not
variables linked by an identity of functions. Thus this calculation does not
prove the claim.
It purports to prove the stronger statement
``$b^n+n\nmid a^n+n$ for $a>b$, with the exception
$1^2+1\mid2^2+2$.'' Can you guess where that exception came from?

\paragraph{What an amazing argument! WAAA!}
For a natural number $n$, prove
\[
\left\lfloor\frac n1\right\rfloor+
\left\lfloor\frac n2\right\rfloor+\cdots+
\left\lfloor\frac nn\right\rfloor+\lfloor\sqrt n\rfloor
\]
is even. (INMO, 2014.)
In $P=\{1,2,\ldots,n\}$, the number of multiples of $q$ is
$\lfloor n/q\rfloor$.
Counting multiples of $k$ not exceeding $n$ is equivalent to counting divisors
of the numbers $1,\ldots,n$:
``Multiples are obtained by multiplying by a constant, and divisors by dividing
by a constant.'' WAAA!
Thus
\[
\sum_{k=1}^n\left\lfloor\frac nk\right\rfloor+\lfloor\sqrt n\rfloor
=\sum_{k=1}^n d(k)+\lfloor\sqrt n\rfloor.
\]
Now $d(k)\equiv1\pmod2$ exactly when $k$ is a perfect square, and $P$
contains $\lfloor\sqrt n\rfloor$ squares. Therefore
\[
\sum_{k=1}^n\left\lfloor\frac nk\right\rfloor+\lfloor\sqrt n\rfloor
\equiv2\lfloor\sqrt n\rfloor\equiv0\pmod2.
\]
\emph{Quite intelligently done.}

% Source page 9 - left


\paragraph{Challenge: Kiran S. Kedlaya, Crux Mathematicorum.}
For positive integers $x,y,z$, prove that $(xy+1)(yz+1)(zx+1)$ is a perfect square
if and only if $xy+1$, $yz+1$, and $zx+1$ are each perfect squares.
Soumik Ghosh communicated a solution that seems deceptively correct.
Hats off to the ingenious use of infinite descent. Dangerous!

Call $(x,y,z)$ good if $(xy+1)(yz+1)(zx+1)$ is a perfect square.
The claim is that every three-element subset of $(x,y,z,t)$ is good, where
\[
t=x+y+z+2xyz-2\sqrt{(xy+1)(yz+1)(zx+1)}.
\]
This $t$ is a root of
\[
x^2+y^2+z^2+t^2-2(xy+yz+zx+xt+yt+zt)-4xyzt-4=0,
\]
which may be rewritten as
\[
(x+y-z-t)^2=4(xy+1)(zt+1),\quad
(y+z-x-t)^2=4(yz+1)(xt+1),\quad
(z+x-y-t)^2=4(zx+1)(yt+1).
\]
Wow! This is equivalent to a Ramanujan move.
Multiply $4(xy+1)(zt+1)$, $4(yz+1)(xt+1)$, and $(xy+1)^2$, then divide by
$(xy+1)(yz+1)(zx+1)$ to obtain that $(xy+1)(yt+1)(zt+1)$ is a perfect square.
Similarly, $(yz+1)(yt+1)(zt+1)$ and $(xz+1)(xt+1)(zt+1)$ are perfect squares.
Thus every three-element subset is good.
Now, without loss of generality, let $z\geq x,y$. Then $z\geq t$, so replace $(z,y,x)$ by $(t,y,x)$.
At every step the maximal element decreases, until eventually $t=0$.
It follows, by arguing correspondingly, that $xy+1$, $yz+1$, and $zx+1$ are perfect squares.
\emph{Quite dangerously done.}

As written, this descent is incomplete: the displayed multiplication does not
directly yield each asserted target product, and the argument has not shown that
$t$ is a nonnegative integer strictly smaller than the element it replaces.

Here is another problem.
\paragraph{Challenge.}
For $a_1,\ldots,a_n,x,y\in\mathbb N$, prove that the following system has solutions if and only if $x=y$:
% Source page 10 - right
\[
a_1^x+\cdots+a_n^x=b^x,\qquad a_1^y+\cdots+a_n^y=b^y
\]
where $b$ is a natural number. As written, $n=1$ and $a_1=b$ is a counterexample unless a nontriviality condition is added.


\subsection{Calculus---oh, not again!}
\paragraph{IMO (1997).}
Solve $a^{b^2}=b^a$ in positive integers.
Let $a\geq b$. The source takes $(a,b)=b$, so $a=bk$.
Then $b^{b^2}k^{b^2}=b^{bk}$ and $k^{b^2}=b^{bk-b^2}$.
Put $k=b^x$. This gives
\[
xb^2=b^{x+1}-b^2
\Longrightarrow b^2(x+1)=b^{x+1}
\Longrightarrow x+1=b^{x-1}.
\]
Set $y=x+1$. Then $y=b^{y-2}$, $k=b^{y-1}$, $a=b^y$, and
$b=\sqrt[y-2]{y}$. The pairs $(b,y)$ are $(1,1),(3,3),(2,4)$.
Since $2^{y-2}-y$ is increasing, the difference grows exponentially.
Thus the only pairs $(a,b)$ are $(1,1),(27,3),(16,2)$.

This proposed reduction has not established that $\gcd(a,b)=b$, that
$k$ is a power $b^x$, or that it is enough to assume $a\geq b$.
Here is a complete argument.
The pair $(1,1)$ is a solution. Otherwise $a$ and $b$ have the same prime divisors.
Comparing prime valuations shows that there are an integer $c>1$ and coprime positive integers $r,s$ such that
$a=c^r$, $b=c^s$, and
\[
r c^{2s}=s c^r.
\]
If $r<2s$, then $r c^{2s-r}=s$; coprimality forces $r=1$, but $s=c^{2s-1}$ is impossible.
The case $r=2s$ is also impossible. Hence $r>2s$, and
$r=s c^{r-2s}$; coprimality now forces $s=1$.
Thus $r=c^{r-2}$ with $r>2$. For $r\geq5$, we have $c^{r-2}\geq2^{r-2}>r$, while
$r=3$ gives $c=3$ and $r=4$ gives $c=2$.
Thus the only pairs $(a,b)$ are $(1,1),(27,3),(16,2)$.

\subsection{Number-theoretic inequalities}
Here is a legendary problem due to Paul Erd\H{o}s.
Each positive integer $a_1,\ldots,a_n$ is less than $1951$, and the least common multiple of any two exceeds $1951$.
Show $1/a_1+\cdots+1/a_n<1+n/1951$.
A counting argument gives the bound.
The number of multiples of $a_i$ in $\{1,\ldots,1951\}$ is $\lfloor1951/a_i\rfloor$, so
$\sum_i\lfloor1951/a_i\rfloor\leq1951$.
Since $x-1<\lfloor x\rfloor$, we obtain
$\sum_i(1951/a_i-1)<1951$, proving the claim.

% Source page 15 - left
Here is another problem of Erd\H{o}s.
\paragraph{Magical solution: Erd\H{o}s--Straus problem.}
Given $3\mid n+1$, prove that there are integers $x,y,z$ with $4xyz=n(xy+yz+zx)$.
Take $x=n$, $y=(n-2)/3+1$, and $z=n(n+1)/3$.
This proves the claim. The case $4\mid n+1$ leads to a related problem.

Two related problems are recorded below.
\paragraph{Problem 1.}
Given $\alpha,\beta\in\mathbb Z^+$ and primes $p_1,\ldots,p_n,q$, show that, for every $\alpha$, there exists an $n$ such that
$p_1^\alpha+p_2^\alpha+\cdots+p_n^\alpha=q^\beta$.
Here $n$ and $\alpha$ are each both fixed and quantified, so the statement needs a corrected choice of quantifiers.
% Source page 15 - right
\paragraph{Problem 2.}
Find infinitely many families of solutions in nine distinct integers
$x_1,x_2,x_3,y_1,y_2,y_3,z_1,z_2,z_3$ of
\[
(x_1^2+x_2^2+x_3^2)+(y_1^2+y_2^2+y_3^2)=z_1^2+z_2^2+z_3^2,
\]
even if no three of the numbers form a Pythagorean triple.
Both admit unexpectedly concise constructions.
% BLOG-CONTENT-END
\end{document}
